\chapter{The Venezuelan decade}
\label{ch:venezuela}

\section{The third patron}

In October 2000 Fidel Castro and Hugo Ch\'avez signed an Integral Cooperation
Agreement in Caracas. Venezuela would supply Cuba with crude oil and refined
products on concessional terms --- initially some 53,000 barrels a day, rising
over the decade to around 90,000 or 100,000 --- and Cuba would pay, in
substantial part, not in money but in the services of its people: doctors,
nurses, dentists, teachers, sports trainers, agricultural technicians and
military and intelligence advisers.\unv{}

It is worth pausing on the structure of this deal, because it is the third time
the same structure appears in this book. In 1934 the United States guaranteed a
price for Cuban sugar above the world market and Cuba organised itself around
it. From 1972 the Soviet Union guaranteed a price for Cuban sugar and supplied
oil below the world market, and Cuba organised itself around that. From 2000
Venezuela supplied oil below the world market and bought Cuban services above
their market value, and Cuba organised itself around that. Each arrangement was
generous, each was political rather than commercial, each rescued the Cuban
economy at the moment it was struck, and each removed the pressure that would
otherwise have forced diversification. Chapter~\ref{ch:sovieteconomy} stated the
rule; this chapter is its third confirmation.

\section{Doctors as an export industry}

The exchange grew far beyond oil. From 2003, under the Barrio Adentro
programme, Cuba placed medical personnel in Venezuelan barrios and rural areas
where Venezuelan doctors had never practised --- at the peak something in the
range of twenty to thirty thousand health workers, with total Cuban civilian
personnel in Venezuela reaching perhaps forty thousand.\unv{} Similar
arrangements, smaller and mostly paid in cash, were signed with Brazil,
Angola, Algeria, South Africa, Qatar and dozens of other countries, and after
2013 the Mais M\'edicos programme placed some 11,000 Cuban doctors in
underserved regions of Brazil.

The economics of this transformed the country. By the early 2010s Cuba's
exports of professional services were worth on the order of eight to ten billion
dollars a year, of which health services were the large majority --- against
tourism at roughly two and a half billion and nickel at about one and a
half \citep{onei,mesalago2013}.\unv{} Cuba's largest export was no longer a commodity. It was its
doctors. That is a genuinely remarkable outcome for a country of eleven
million, it is the direct dividend of the investment described in
Chapter~\ref{ch:socialstate}, and it is the strongest existing evidence for the
proposition in Chapter~\ref{ch:talent} that Cuba's educated population is an
exportable asset.

The arrangement's other half has to be reported as plainly. The host government
paid Cuba; Cuba paid the doctor. The doctor's share was typically a small
fraction of the contract value --- in the Brazilian programme, the most
transparently documented, Cuba received roughly four times what it passed to the
physician \citep{kirk2015}.\unv{} Doctors on mission were subject to rules that would not be
accepted in most labour markets: in several destinations their passports were
held by the mission chief; movement, association and relationships with locals
were restricted; and a doctor who abandoned the mission was, until the
migration reform of 2013, barred from returning to Cuba for eight years, which
in practice meant separation from children. The United States operated a Cuban
Medical Professional Parole programme from 2006 to 2017 specifically to induce
defections, and several thousand took it.

Both descriptions are true at once, and the honest summary is that Cuba built a
real and valuable export industry and financed it partly by underpaying and
over-controlling the people who produced it. Chapter~\ref{ch:life} argues that
the programme should be kept and the labour arrangement changed, and that this
is not only the ethical answer but the revenue-maximising one, since the
industry's binding constraint is now the supply of willing doctors.

\section{What the money bought}

The Venezuelan relationship coincided with, and paid for, the most politically
confident period the Cuban government has had since the 1980s.

Abroad, it bought a bloc. ALBA was founded by Cuba and Venezuela in 2004 and
grew to include Bolivia, Nicaragua, Ecuador and several Caribbean states.
Petrocaribe, from 2005, extended Venezuelan concessional oil across the
Caribbean with Cuban medical missions attached. The regional isolation that had
been American policy since 1962 simply dissolved: the Organization of American
States revoked Cuba's 1962 suspension in 2009, the Community of Latin American
and Caribbean States was founded without the United States in 2011 and Cuba
chaired it in 2013, and the European Union's 1996 Common Position, which had
conditioned relations on democratisation, was steadily bypassed by member states
and formally replaced in 2016.

At home it bought the Battle of Ideas. This began in late 1999 with the case of
Eli\'an Gonz\'alez, a five-year-old boy who survived the sinking of his mother's
boat and was held in Miami by relatives against his father's wishes. Cuba
mobilised around his return for seven months, in daily marches and a permanent
tribune built opposite the American mission, and won: the boy went home in June
2000. It was the last unambiguous political victory the Cuban government won on
its own terms, and the mobilisation machinery it built was then turned to
domestic programmes --- an army of young social workers, university courses
delivered in every municipality, television as a teaching medium, and a great
deal of capital spending decided by mass campaign rather than by plan.

One of those campaigns deserves its own paragraph, because Part III depends on
it. By 2005 the Cuban grid was failing: the thermal plants were beyond their
design lives, unplanned outages were routine, and a hurricane season could take
the country dark. The Energy Revolution of 2006 attacked the problem from both
ends at once. On the demand side the state replaced essentially every
incandescent bulb in the country --- some nine million --- with compact
fluorescents, distributed millions of efficient rice cookers and pressure
cookers, and replaced the oldest refrigerators and air conditioners. On the
supply side it installed around four thousand small distributed generator sets,
roughly 1,300 megawatts in total, sited across the country rather than
concentrated in a few large plants.\unv{} Blackout hours fell sharply. The
programme was expensive, the generators burned expensive fuel, and the strategy
was later allowed to decay --- but the diagnosis was right, and it was right in
a way that matters now: Cuba's grid problem is best attacked by distributed
capacity and demand reduction, not by rebuilding large central plants.
Chapter~\ref{ch:energy} builds directly on this precedent, with the generators
replaced by photovoltaics and storage.

\section{Two shocks, and a credit history}

In August, September and November of 2008 three hurricanes --- Gustav, Ike and
Paloma --- crossed Cuba within ten weeks. Together they damaged or destroyed
something over half a million homes and caused losses conventionally estimated
at around ten billion dollars, on the order of a fifth of national
output.\unv{} It was the worst hurricane season in Cuban recorded history, and
it arrived in the same quarter as the global financial crisis.

Cuba's response to the resulting liquidity squeeze did lasting damage. Beginning
in late 2008 the state froze the bank accounts of foreign companies operating in
Cuba --- money those companies had earned and were entitled to repatriate,
amounting to something over a billion dollars --- and simply stopped paying
suppliers.\unv{} Some balances were released over the following years; some were
not. The measure solved a cash problem and created a reputation problem, and
the reputation problem has proved much more durable. Foreign companies now price
Cuban counterparty risk on the basis of what happened in 2009, and no statement
of intent from Havana has changed that. Chapter~\ref{ch:capital} argues that
this is a credit history, that credit histories are repaired by behaviour over
years rather than by legislation, and that any investment strategy which does
not begin by settling these arrears is proposing to build on a default.

\section{The end of it}

Hugo Ch\'avez was diagnosed with cancer in 2011 and was treated in Havana. He
died on 5 March 2013. Venezuela's own economy then entered the collapse that has
defined it since, and the terms Cuba had enjoyed began to erode --- shipments
falling year by year through the mid-2010s, and falling faster after 2016.

By then Cuba had had thirteen years of subsidised oil and a services export
boom, and had used them, as it had used the Soviet years, to postpone rather
than to restructure. The sugar industry was not rebuilt. Domestic food
production did not rise. The dual currency was not unified, although unification
was announced as policy in 2011 and prepared for a decade. The energy system was
not converted. What the country had instead was one more external rent, larger
than it looked, and a decade in which the pressure to change had been removed.

\begin{ledger}{the Venezuelan decade, 2000--2013}
\emph{Achieved.} An export industry built out of the education system that
displaced sugar, nickel and tourism combined, and that placed Cuban doctors in
places no other country would staff --- a real service really delivered, and
the best evidence available for the argument of Part~III. The end of Cuba's
regional isolation, with the OAS suspension revoked and CELAC founded and
chaired. The Energy Revolution of 2006, a correct diagnosis and a partially
correct cure, whose logic Chapter~\ref{ch:energy} adopts. Recovery from the
worst hurricane season on record without mass casualties, which is what a
civil-defence system with block-level reach is for.

\emph{Cost.} The third guaranteed buyer at a political price in seventy years,
taken for the third time, with the third identical result: no diversification,
no restructuring, and a cliff at the end of it. An export industry financed by
paying its workers a fraction of their contract value and, in several
destinations, holding their passports. The freezing of foreign companies'
accounts in 2009, which bought a year of liquidity at the price of a credit
history that is still being paid for. And thirteen years in which every
structural problem this book has named was identified, announced as policy, and
not fixed.
\end{ledger}
